% ============================================================================
% Section III: Mathematical Formulation
% ============================================================================

\section{Mathematical Formulation}
\label{sec:formulation}

\subsection{Sets and Indices}

The GTEP problem is defined over the following sets:

\begin{itemize}
  \item $\setS$: scenarios, indexed by $s$, with probability $\pi_s$
  \item $\setT$: stages (planning periods), indexed by $t$
  \item $\setB_t$: blocks (operating periods) in stage $t$, indexed by $b$,
        with duration $\Delta_b$ [h]
  \item $\setN$: buses (electrical nodes), indexed by $n$
  \item $\setG$: generators, indexed by $g$, with bus $n_g$
  \item $\setD$: demands, indexed by $d$, with bus $n_d$
  \item $\setL$: transmission lines, indexed by $\ell$, connecting buses
        $a_\ell$ and $b_\ell$
  \item $\setE$: batteries (energy storage), indexed by $e$
  \item $\mathcal{V}$: converters (battery$\leftrightarrow$grid), indexed by $v$
  \item $\setJ$: hydraulic junctions, indexed by $j$
  \item $\setR$: reservoirs, indexed by $r$, with junction $j_r$
  \item $\setW$: waterways, indexed by $w$
  \item $\setH$: turbines, indexed by $h$
  \item $\setZ$: reserve zones, indexed by $z$
  \item $\mathcal{P}$: reserve provisions, indexed by $p$
  \item $\setFR$: flow rights (flow-based water rights), indexed by $i$
  \item $\setVR$: volume rights (volume-based water rights), indexed by $o$
  \item $\setHP$: hydraulic pumps, indexed by $h$
  \item $\setLNG$: LNG storage terminals, indexed by $\tau$
  \item $\setUC$: generators with unit commitment, indexed by $g$
  \item $\setIZ$: inertia zones, indexed by $z$
  \item $\setIP$: inertia provisions, indexed by $p$
\end{itemize}

\begin{figure}[t]
  \centering
  \begin{tikzpicture}[
    level 1/.style={sibling distance=22mm, level distance=14mm},
    level 2/.style={sibling distance=11mm, level distance=14mm},
    scen/.style={draw, rounded corners=3pt, fill=blue!12,
      font=\footnotesize, minimum width=14mm, minimum height=6mm},
    stg/.style={draw, rounded corners=2pt, fill=orange!15,
      font=\footnotesize, minimum width=12mm, minimum height=6mm},
    blk/.style={draw, fill=gray!15,
      font=\scriptsize, minimum width=7mm, minimum height=5mm},
    edge from parent/.style={draw, -{Stealth[length=3pt]}, thick},
    root/.style={draw, rounded corners=4pt, fill=black!8,
      font=\footnotesize\bfseries, minimum width=18mm,
      minimum height=7mm},
  ]
    \node[root] {Problem}
      child { node[scen] {$s_1\;(\pi_1)$}
        child { node[stg] {$t_1\;(\delta_1)$}
          child { node[blk] {$b_1$} }
          child { node[blk] {$b_2$} }
        }
        child { node[stg] {$t_2\;(\delta_2)$}
          child { node[blk] {$b_1$} }
          child { node[blk] {$b_2$} }
          child { node[blk] {$b_3$} }
        }
      }
      child { node[scen] {$s_2\;(\pi_2)$}
        child { node[stg] {$t_1$}
          child { node[blk] {$b_1$} }
          child { node[blk] {$b_2$} }
        }
        child { node[stg] {$t_2$}
          child { node[blk] {$b_1$} }
          child { node[blk] {$b_2$} }
          child { node[blk] {$b_3$} }
        }
      };
    % Level labels on the right
    \node[anchor=west, font=\scriptsize\itshape, text=black!50]
      at (2.9,0) {Scenarios $\setS$};
    \node[anchor=west, font=\scriptsize\itshape, text=black!50]
      at (2.9,-1.4) {Stages $\setT$};
    \node[anchor=west, font=\scriptsize\itshape, text=black!50]
      at (2.9,-2.8) {Blocks $\setB_t\;(\Delta_b)$};
  \end{tikzpicture}
  \caption{Three-level time hierarchy.  The problem decomposes into
    scenarios~$s$ (probability~$\pi_s$), stages~$t$ (discount
    factor~$\delta_t$), and blocks~$b$ (duration~$\Delta_b$).}
  \label{fig:time_hierarchy}
\end{figure}

\subsection{Decision Variables}

\noindent\textbf{Operational variables} (per scenario $s$, stage $t$, block $b$):
\begin{align}
  p_{g,s,t,b}  &\in [\underline{p}_g, \overline{p}_g]
    & \text{generator output [MW]} \label{eq:pgen} \\
  \ell_{d,s,t,b} &\in [0, \overline{\ell}_d]
    & \text{served demand [MW]} \label{eq:load} \\
  q_{d,s,t,b}  &\geq 0
    & \text{unserved demand [MW]} \label{eq:fail} \\
  f_{\ell,s,t,b} &\in [-\overline{F}^{ba}_\ell, \overline{F}^{ab}_\ell]
    & \text{line flow [MW]} \label{eq:flow} \\
  \theta_{n,s,t,b} &\in \R
    & \text{voltage angle [rad]} \label{eq:theta}
\end{align}

\noindent\textbf{Storage state variables}:
\begin{align}
  e_{e,s,t,b}  &\in [\underline{e}_e, \overline{e}_e]
    & \text{battery SoC [MWh]} \label{eq:soc} \\
  v_{r,s,t,b}  &\in [\underline{v}_r, \overline{v}_r]
    & \text{reservoir volume [hm\textsuperscript{3}]} \label{eq:vol} \\
  v_{\tau,s,t,b}^{\text{lng}}
               &\in [\underline{V}_\tau, \overline{V}_\tau]
    & \text{LNG tank volume [m\textsuperscript{3}]} \label{eq:lng_vol}
\end{align}

\noindent\textbf{Unit commitment variables} (per scenario $s$, stage $t$, period $p$):
\begin{align}
  u_{g,s,t,p} &\in \{0,1\} & \text{unit status (1 = online)} \\
  v_{g,s,t,p} &\in \{0,1\} & \text{startup event} \\
  w_{g,s,t,p} &\in \{0,1\} & \text{shutdown event}
\end{align}

\noindent\textbf{Investment variables} (per stage $t$):
\begin{equation}
  m_{g,t} \in [0, \overline{m}_g]
  \quad \text{expansion modules built} \label{eq:expmod}
\end{equation}

\subsection{Objective Function}

Minimize total expected discounted cost:
\begin{equation}
\label{eq:objective}
\min \sum_{s \in \setS} \pi_s \sum_{t \in \setT} \delta_t
\left[
  \sum_{b \in \setB_t} \Delta_b \cdot C^\text{op}_{s,t,b}
  + C^\text{inv}_t
\right]
\end{equation}
where $\delta_t = (1+r)^{-\tau_t/8760}$ is the discount factor with annual
rate $r$ and $\tau_t$ the cumulative hours to stage $t$.

The operational cost for each $(s,t,b)$ is:
\begin{multline}
\label{eq:opcost}
C^\text{op}_{s,t,b} = \sum_{g \in \setG} c_g \, p_{g,s,t,b}
  + \sum_{d \in \setD} c^f_d \, q_{d,s,t,b} \\
  + \sum_{\ell \in \setL} c_\ell
    (f^+_{\ell,s,t,b} + f^-_{\ell,s,t,b})
\end{multline}

The investment cost for stage $t$ is:
\begin{multline}
\label{eq:invcost}
C^\text{inv}_t = \sum_{g \in \setG} \omega^\text{cap}_{g,t} m_{g,t}
  + \sum_{d \in \setD} \omega^\text{cap}_{d,t} m_{d,t} \\
  + \sum_{\ell \in \setL} \omega^\text{cap}_{\ell,t} m_{\ell,t}
  + \sum_{e \in \setE} \omega^\text{cap}_{e,t} m_{e,t}
\end{multline}
where $\omega^\text{cap}_{(\cdot),t} = K_{(\cdot)} \delta_t T_t /
\sigma$ is the weighted investment cost coefficient,
$c_g$ is the variable generation cost [\$/MWh], $c^f_d$ is the
demand failure penalty [\$/MWh], $c_\ell$ is the transmission cost
[\$/MWh], $K_{(\cdot)}$ is the annualized investment cost
[\$/year] per module for each component, and $T_t$ is the stage
duration [h].  All objective coefficients are divided by a
scaling factor $\sigma$ (\texttt{scale\_objective}, default 1000) to
improve solver numerical conditioning.  Concretely, each operational LP
cost coefficient is weighted by $\omega_{s,t,b} = \pi_s \delta_t
\Delta_b / \sigma$, so the reported objective equals $z/\sigma$.

\subsection{Bus Power Balance}
\label{sec:balance}
\begin{multline}
\label{eq:balance}
\sum_{\substack{g \in \setG \\ n_g = n}} (1{-}\lambda_g) p_{g,s,t,b}
- \sum_{\substack{d \in \setD \\ n_d = n}}
  (1{+}\lambda_d)\,\ell_{d,s,t,b} \\
+ \sum_{\ell \in \setL} A_{n,\ell} \, f_{\ell,s,t,b} = 0
\end{multline}
where $\lambda_g, \lambda_d$ are injection and withdrawal loss factors,
respectively.  The line flow terms use directional contributions rather
than a simple incidence matrix: for each line~$\ell$ from bus~$a_\ell$
to~$b_\ell$, the sending bus sees $-f^+_\ell$ (or $+(1-\lambda_\ell)
f^-_\ell$ when losses are modeled) and the receiving bus sees
$+(1-\lambda_\ell)f^+_\ell$ (or $-f^-_\ell$), where $\lambda_\ell$
is the line loss factor.  When $\lambda_\ell=0$, a single bidirectional
variable suffices and the incidence reduces to $A_{n,\ell} \in
\{+1,-1,0\}$.  The dual of~\eqref{eq:balance} is the locational marginal
price (LMP) at bus~$n$.

\subsection{Demand Balance}

When demand curtailment is allowed, served and unserved demand sum to the
scheduled load:
\begin{equation}
\label{eq:demand_balance}
\ell_{d,s,t,b} + q_{d,s,t,b} = \overline{\ell}_{d,t,b}
\quad \forall\; d, s, t, b
\end{equation}
where $\overline{\ell}_{d,t,b}$ is the scheduled demand at block~$b$,
incorporating any demand profile scaling.  When curtailment is disabled,
the load is fixed: $\ell_{d,s,t,b} = \overline{\ell}_{d,t,b}$.

Instead of the single penalty $c^f_d$, a demand can carry a stepwise
value-of-lost-load curve of $K$ segments (\texttt{fail\_segments}): the
served load is split into $K$ tranches bounded by
$f_{d,k}\,\overline{\ell}_{d,t,b}$, and the curtailment of tranche~$k$
is priced at $c^f_{d,k}$.  With non-decreasing segment costs the LP
curtails the cheapest segment first, so the curve needs no extra rows,
ordering constraints or binaries.

\subsection{DC Power Flow (Kirchhoff Constraints)}

When \texttt{use\_kirchhoff} is enabled, for each line $\ell$ and
$(s,t,b)$:
\begin{equation}
\label{eq:kirchhoff}
f_{\ell,s,t,b} = \frac{V_\ell^2}{X_\ell}
  \left(\theta_{a_\ell,s,t,b} - \theta_{b_\ell,s,t,b}\right)
\end{equation}
where $X_\ell$ is the line reactance [$\Omega$] and $V_\ell$ is the
nominal voltage [kV].  The angle variables are enlarged by a factor
$\sigma_\theta$ (\texttt{scale\_theta}) for numerical conditioning:
$\theta^\text{LP} = \sigma_\theta \cdot \theta$.  When it is not set,
$\sigma_\theta = 1 / \bigl(3 \operatorname{median}_\ell (X_\ell /
V_\ell^2)\bigr)$.  An angle accumulates about $D$ line drops on the way
to the reference bus, and the factor $3 \approx \sqrt{D}$ balances the
angle coefficient of the Kirchhoff row against the size of
$\theta^\text{LP}$.  A value set by the user is kept as given.

A line with zero reactance, or with $|X_\ell/V_\ell^2|$ below a
threshold (default $10^{-6}$~p.u.), is a zero-impedance tie by default:
the flow term of~\eqref{eq:transformer} vanishes, the two bus angles are
forced equal (up to the phase shift), and the flow is free within the
line limits.  A line without a reactance is a DC link and has no
Kirchhoff row.

Line flow is bounded by transfer capacities:
\begin{equation}
-\overline{F}^{ba}_\ell \leq f_{\ell,s,t,b}
  \leq \overline{F}^{ab}_\ell
\end{equation}

\paragraph{Loop-flow formulation and switching.}
Besides the angle-based form~\eqref{eq:kirchhoff}, Kirchhoff's voltage
law can be written over a cycle basis of the network (one row per
independent loop, with no angle variables), which is smaller for
radial-heavy systems.  Lines can be declared switching candidates with a
binary state; the Kirchhoff row of a switchable line (or of every cycle
through it) is relaxed with a big-M disjunction, giving optimal
transmission switching as a MIP.

\paragraph{Transformer model.}
For tap-changing and phase-shifting transformers, the Kirchhoff
constraint generalizes to:
\begin{equation}
\label{eq:transformer}
-\theta_{a_\ell} + \theta_{b_\ell}
  + (\tau_\ell \chi_\ell)(f^+_\ell - f^-_\ell)
  = -\sigma_\theta \, \psi_{\ell}
\end{equation}
where $\tau_\ell$ is the off-nominal tap ratio (default~1),
$\chi_\ell = \sigma_\theta X_\ell / V_\ell^2$ is the scaled reactance,
and $\psi_\ell$ is the phase-shift angle [rad].
When $\tau_\ell = 1$ and $\psi_\ell = 0$,~\eqref{eq:transformer}
reduces to~\eqref{eq:kirchhoff}.

\subsection{Line Losses}
\label{sec:losses}

Losses are modeled with one of several formulations, selected per line
or system-wide.

\paragraph{Linear and piecewise models.}
The linear model splits the flow into forward and reverse components
$f^+_\ell, f^-_\ell \geq 0$; the receiving bus sees
$(1-\lambda_\ell)$ of the power sent.  The piecewise model approximates
the quadratic loss $R_\ell f^2 / V_\ell^2$ with $K$ equal segments per
direction, so the loss slope increases with the flow.

\paragraph{Loss arbitrage.}
Both models are exact only when the LP has no incentive to lose
energy.  When the sum of the duals of the two end buses is negative,
$\pi_{a_\ell} + \pi_{b_\ell} < 0$, which happens with curtailed
renewables, inflexible units or congestion with angular coupling, the
LP can fill expensive segments first or circulate flow in both
directions at once.  The losses then become a fictitious energy sink
that sets prices at zero instead of letting them become negative.
Section~\ref{sec:cases} shows the size of this effect on the Chilean
system.

\paragraph{Tangent outer approximation with a loss price.}
The \texttt{tangent\_signed\_flow} model uses a signed flow
$f_\ell$ and a loss variable $\lambda^{\mathrm{L}}_\ell$ bounded below
by $K$ tangents of the convex quadratic~\cite{coffrin2014lpac},
\begin{equation}
\label{eq:loss_tangent}
\lambda^{\mathrm{L}}_{\ell,s,t,b} \geq 2 c_\ell f^{(k)}_\ell
  f_{\ell,s,t,b} - c_\ell \bigl(f^{(k)}_\ell\bigr)^2,
  \quad k = 1,\dots,K,
\end{equation}
with $c_\ell = R_\ell/V_\ell^2$.  It is bounded above by a secant chord
on an auxiliary $v_\ell \geq |f_\ell|$, which closes the
bidirectional-flow channel.  A small loss price $\varepsilon_\ell$
[\$/MWh] on $\lambda^{\mathrm{L}}_\ell$, weighted by
$\omega_{s,t,b}$, removes the remaining incentive to inflate the losses
whenever $\varepsilon_\ell \geq \tfrac12 |\pi_{a_\ell}+\pi_{b_\ell}|$.
Its effect on legitimate operation is proportional to the losses, not to
the flow.  The same price is applied to the loss columns of the
piecewise and bidirectional models, where it also breaks the degeneracy
between one-way and two-way flow.  It defaults to 0.1~\$/MWh, can be
raised per line on the corridors where the arbitrage appears, and is
disabled with $\varepsilon_\ell = 0$.
For lines where the chord is too loose, an $L$-segment chord
with a type-2 special ordered set makes the upper bound tight at the
breakpoints, with a worst-case error of
$c_\ell \overline{F}_\ell^2/(4L^2)$.

\subsection{Battery Energy Storage}
\label{sec:battery}

State-of-charge balance for battery $e$ between consecutive blocks:
\begin{multline}
\label{eq:soc_balance}
e_{e,s,t,b} = e_{e,s,t,b-1}(1 - \mu_e^h \Delta_b) \\
  + p^\text{in}_{e,s,t,b} \cdot \eta^\text{in}_e \cdot \Delta_b
  - p^\text{out}_{e,s,t,b} \cdot \Delta_b / \eta^\text{out}_e
\end{multline}
where $\eta^\text{in}_e$ and $\eta^\text{out}_e$ are charge and discharge
efficiencies, $p^\text{in}, p^\text{out}$ are charge and discharge
power [MW], and $\mu_e^h = \mu_e / 8760$ is the hourly self-discharge
rate derived from the annual standing-loss rate~$\mu_e$
(\texttt{annual\_standing\_loss}, default 0.48~p.u./year, about 4\% per
month).
Initial SoC $e_{e,\cdot,\cdot,0} = e^\text{ini}_e$ is an
equality constraint in the first stage; the final SoC of stage~$t$
links to the initial SoC of stage~$t{+}1$.

\begin{figure}[t]
  \centering
  \resizebox{\columnwidth}{!}{%
  \begin{tikzpicture}[
    >=Stealth,
    axis/.style={->, thick},
    charge/.style={fill=green!15},
    discharge/.style={fill=red!10},
  ]
    % Axes
    \draw[axis] (0,0) -- (7.2,0)
      node[right, font=\scriptsize] {Block (hour)};
    \draw[axis] (0,0) -- (0,3.2)
      node[above, font=\scriptsize] {SoC (\%)};

    % Y-axis ticks
    \foreach \y/\lbl in {0/0, 0.6/20, 1.2/40, 1.8/60, 2.4/80, 3.0/100} {
      \draw (-0.08,\y) -- (0.08,\y)
        node[left=2pt, font=\tiny] {\lbl};
    }
    % X-axis ticks
    \foreach \x/\lbl in
      {0.3/1, 1.5/6, 2.7/10, 3.6/13, 4.5/16, 5.7/20, 6.9/24} {
      \draw (\x,-0.08) -- (\x,0.08)
        node[below=2pt, font=\tiny] {\lbl};
    }

    % SoC data points
    \coordinate (p0) at (0.0, 0.9);   % hour 0:  30%
    \coordinate (p1) at (0.9, 0.6);   % hour 3:  20%
    \coordinate (p2) at (1.8, 0.6);   % hour 6:  20%
    \coordinate (p3) at (2.7, 1.5);   % hour 10: 50%
    \coordinate (p4) at (3.6, 2.7);   % hour 13: 90%
    \coordinate (p5) at (4.2, 2.7);   % hour 15: 90%
    \coordinate (p6) at (5.1, 1.5);   % hour 18: 50%
    \coordinate (p7) at (6.0, 0.9);   % hour 21: 30%
    \coordinate (p8) at (6.9, 0.9);   % hour 24: 30%

    % Charge shading (hours 6--15)
    \fill[charge] (p2) -- (p3) -- (p4) -- (p5)
      -- (p5 |- 0,0) -- (p2 |- 0,0) -- cycle;

    % Discharge shading (hours 15--21)
    \fill[discharge] (p5) -- (p6) -- (p7)
      -- (p7 |- 0,0) -- (p5 |- 0,0) -- cycle;

    % SoC line
    \draw[thick, blue!80] (p0) -- (p1) -- (p2) -- (p3) -- (p4)
      -- (p5) -- (p6) -- (p7) -- (p8);

    % Labels
    \node[font=\tiny, green!50!black, above] at (3.15,2.1)
      {charge};
    \node[font=\tiny, red!60!black, above] at (4.95,1.8)
      {discharge};

    % Solar peak annotation
    \draw[{Stealth[length=3pt]}-, thin, gray]
      (3.6,2.7) -- (4.3,3.1)
      node[right, font=\tiny] {solar peak};

    % Evening peak annotation
    \draw[{Stealth[length=3pt]}-, thin, gray]
      (5.1,1.5) -- (5.8,2.2)
      node[right, font=\tiny] {evening peak};
  \end{tikzpicture}}%
  \caption{Battery SoC over 24~blocks: charge during solar hours (green),
    discharge during evening peak (red).}
  \label{fig:battery_soc}
\end{figure}

\subsection{Converter Constraints}

A converter~$v$ couples a battery~$e$ to the electrical network through
a generator~$g$ (discharge path) and a demand~$d$ (charge path) via a
power conversion rate~$\rho_v$:
\begin{align}
\label{eq:conv_discharge}
p_{g,s,t,b} &= \rho_v \cdot p^\text{out}_{e,s,t,b} \\
\label{eq:conv_charge}
\ell_{d,s,t,b} &= \rho_v \cdot p^\text{in}_{e,s,t,b}
\end{align}
When the converter has a capacity limit~$\overline{C}_v$, a joint bound
applies: $p_{g,s,t,b} + \ell_{d,s,t,b} \leq \bar{C}_{v,t}$.
Converter capacity is expandable using the same modular framework
as~\eqref{eq:capacity}.

\subsection{Hydro Cascade}
\label{sec:hydro}

For each junction $j$ with reservoir $r$, the water balance is:
\begin{multline}
\label{eq:hydro_balance}
v_{r,s,t,b} = v_{r,s,t,b-1} \\
  + \alpha_r \Delta_b \!\Bigl(
    \textstyle\sum_{w \in \text{in}(j)} q_{w}
    - \textstyle\sum_{w \in \text{out}(j)} q_{w}
    + F_j
  \Bigr)
\end{multline}
where $q_w$ is the waterway discharge [m$^3$/s], $F_j$ is the exogenous
inflow at junction $j$, and $\alpha_r = 0.0036$ converts flow rates to
volume increments: since 1~m$^3$/s flowing for 1~hour equals
3600~m$^3$ = 0.0036~hm$^3$, $\alpha_r = 3.6 \times 10^{-3}$
[hm$^3$/(m$^3$/s$\cdot$h)] (the default when \texttt{flow\_conversion\_rate}
is not set).  A junction can release surplus water through a drain
column bounded by a capacity and priced, like waterway flows and
spillways, per (m$^3$/s)$\cdot$h with the block weight
$\omega_{s,t,b}$.  These exit costs are optional.  When every exit of a
full reservoir carries a cost~$c$, its water value (the dual
of~\eqref{eq:hydro_balance}) is $-c/\alpha_r$: one more hm$^3$ must be
discharged at a cost, and SDDP carries such negative values through its
cuts.

Each turbine $h$ linking waterway $w$ to generator $g$ provides:
\begin{equation}
\label{eq:turbine}
p_{g,s,t,b} \leq \rho_h \cdot q_{w,s,t,b}
\end{equation}
where $\rho_h$ is the water-to-power conversion rate [MW$\cdot$s/m$^3$].

\paragraph{Volume-dependent turbine efficiency.}
In practice, the conversion rate depends on the hydraulic head, which
varies with reservoir volume.  When defined, $\rho_h$ is replaced by a
concave piecewise-linear function:
\begin{equation}
\label{eq:vol_efficiency}
\rho_h(v) = \min_{i=1}^{N}
  \bigl\{ c_i + m_i (v - v_i) \bigr\}
\end{equation}
where $(v_i, m_i, c_i)$ are volume breakpoints, slopes, and intercepts
of~$N$ segments.  In the SDDP method, the active segment is updated at
each forward-pass iteration using the current reservoir volume.

\paragraph{Filtration (seepage).}
Seepage from a waterway into a reservoir is modeled as a linear function
of the average volume:
\begin{equation}
\label{eq:filtration}
\varphi_{i,s,t,b} = a_i + b_i \cdot
  \frac{v_{r,s,t,b-1} + v_{r,s,t,b}}{2}
\end{equation}
where $a_i$ is a constant term and $b_i$ is the seepage rate per unit
volume.  As with turbine efficiency, this admits a piecewise-linear
extension selected by the last known volume.

\begin{figure}[t]
  \centering
  \begin{tikzpicture}[
    res/.style={trapezium, trapezium angle=75, draw, fill=cyan!15,
      minimum width=1.4cm, minimum height=0.7cm, font=\scriptsize,
      trapezium stretches body},
    junc/.style={circle, draw, fill=gray!20, minimum size=0.4cm,
      font=\tiny, inner sep=1pt},
    turb/.style={regular polygon, regular polygon sides=3, draw,
      fill=yellow!25, minimum size=0.5cm, font=\tiny, inner sep=0pt},
    gen/.style={circle, draw, fill=orange!20, minimum size=0.5cm,
      font=\tiny},
    flow/.style={-{Stealth[length=4pt]}, thick, blue!70},
    spill/.style={-{Stealth[length=4pt]}, thick, red!50, dashed},
  ]
    % Reservoir 1
    \node[res] (r1) at (0,2.2) {$R_1$};
    \node[font=\tiny\itshape, text=black!70, right=1pt of r1] {$v_1$};

    % Junction 1
    \node[junc] (j1) at (0,1.1) {$j_1$};

    % Inflow
    \draw[flow] (-1.2,2.8) -- (j1.north west)
      node[midway, left, font=\tiny] {inflow $F_1$};

    % R1 -> J1
    \draw[flow] (r1) -- (j1);

    % Turbine
    \node[turb] (t1) at (1.2,0.6) {};
    \node[font=\tiny, right=1pt of t1] {$h_1$};

    % Junction 2
    \node[junc] (j2) at (0,-0.2) {$j_2$};

    % J1 -> T1
    \draw[flow] (j1) -- (t1) node[midway, above, font=\tiny] {$q_w$};

    % T1 -> Generator
    \node[gen] (g1) at (2.5,0.6) {\raisebox{-0.5pt}{\scalebox{0.5}{$\sim$}}};
    \node[font=\tiny, right=1pt of g1] {gen};
    \draw[-{Stealth[length=4pt]}, thick, orange!70] (t1) -- (g1);

    % J1 -> J2 spillway
    \draw[spill] (j1) -- (j2) node[midway, left, font=\tiny] {spill};

    % Reservoir 2
    \node[res] (r2) at (2,-0.2) {$R_2$};
    \node[font=\tiny\itshape, text=black!70, right=1pt of r2] {$v_2$};

    % Junction 3
    \node[junc] (j3) at (1,-1.2) {$j_3$};

    % J2 -> R2 waterway
    \draw[flow] (j2) -- (r2);

    % T1 -> J3 (tailwater)
    \draw[flow] (t1) |- (j3);

    % R2 -> J3
    \draw[flow] (r2) -- (j3);

    % Outflow
    \draw[flow] (j3) -- (1,-2.0) node[below, font=\tiny] {outflow};
  \end{tikzpicture}
  \caption{Hydro cascade example: two reservoirs ($R_1$, $R_2$),
    three junctions ($j_1$--$j_3$), one turbine ($h_1$) driving a
    generator, with inflow, spillway, and waterway connections.}
  \label{fig:hydro_cascade}
\end{figure}

\subsection{Capacity Expansion}

Installed capacity evolves as:
\begin{equation}
\label{eq:capacity}
\bar{C}_{g,t} = \bar{C}_{g,t-1}(1-\xi_{g,t}) + M_g \cdot m_{g,t}
\end{equation}
where $\bar{C}_{g,0}$ is the initial installed capacity, $\xi_{g,t}$ is
the stage derating factor, $M_g$ is the capacity per expansion module
(\texttt{expcap}), and $m_{g,t}$ is the number of modules built.  The
generator output is bounded by the installed capacity:
\begin{equation}
p_{g,s,t,b} \leq \bar{C}_{g,t}
\end{equation}
Analogous expansion constraints apply to demands, lines, batteries, and
converters.

\subsection{Generator Profile Constraints}

When a generator has a capacity profile $\phi_g(s,t,b)$ (e.g., solar
irradiance or wind availability), the profile bounds the output.  The
profile is a per-unit factor of the capacity, and the available power is
capped by the generation range $\overline{p}_{g,t,b}$:
\begin{equation}
\label{eq:profile}
p_{g,s,t,b} \leq A_{g,s,t,b}
  = \min\bigl(\phi_g(s,t,b) \cdot \bar{C}_{g,t},\;
    \overline{p}_{g,t,b}\bigr)
\end{equation}
When the capacity is an expansion variable the row reads
$p_{g,s,t,b} - \min(\phi_g, 1)\,\bar{C}_{g,t} \leq 0$.  Curtailment
needs no variable of its own: the spilled renewable energy is computed
after the solve as $s_{g,s,t,b} = \max(0, A_{g,s,t,b} - p_{g,s,t,b})$,
and its marginal value is the dual of~\eqref{eq:profile}.  Only when
curtailment carries a cost $c_g^\text{spill}$ is the equality
$p_{g,s,t,b} + s_{g,s,t,b} = A_{g,s,t,b}$ with
$0 \leq s_{g,s,t,b} \leq A_{g,s,t,b}$ and cost
$c_g^\text{spill}\,\omega_{s,t,b}\,s_{g,s,t,b}$ used instead.

\subsection{Demand Minimum Energy}

An optional minimum energy constraint ensures that a demand~$d$ is
served a minimum total energy within each stage:
\begin{equation}
\label{eq:min_energy}
\sum_{b \in \setB_t} \ell^\text{man}_{d,s,t,b} \cdot \Delta_b
  \geq E_d^\text{min}
\end{equation}
where $\ell^\text{man}_{d,s,t,b}$ is the mandatory served load and
$E_d^\text{min}$ [MWh] is the stage energy requirement.

\subsection{Reserve Constraints}
\label{sec:reserve}

For each reserve zone $z$ and $(s,t,b)$:
\begin{equation}
\label{eq:reserve_up}
\sum_{p \in \mathcal{P}_z} \gamma_{p,t} \, r^\text{up}_{p,s,t,b}
  + q^\text{up}_{z,s,t,b}
  \geq R^\text{up}_{z,s,t,b}
\end{equation}
where $r^\text{up}_{p,s,t,b}$ is the up-reserve provision from
provision~$p$, $\gamma_{p,t}$ is the provision factor (typically~1),
$q^\text{up}_{z,s,t,b}$ is unserved reserve, and
$R^\text{up}_{z,s,t,b}$ is the zone requirement.  A symmetric constraint
exists for down-reserve.  Reserve provision is bounded by the generator's
headroom: $p_{g,s,t,b} + r^\text{up}_{p,s,t,b} \leq \bar{C}_{g,t}$.

\subsection{Hydro Pump}
\label{sec:pump}

A \texttt{Pump} element $h \in \setHP$ models a hydraulic pump
that consumes electrical power from a demand $d$ and raises water
through a waterway $w$ from a downstream junction to an upstream
junction---the reverse of a turbine.  The conversion constraint ties
the electrical load to the pumped flow:
\begin{equation}
\label{eq:pump}
\frac{\phi_{h,t}}{\eta_{h,t}} \, \varphi_{w,s,t,b} \leq \ell_{d,s,t,b}
\qquad \forall\, h \in \setHP,\; s,t,b
\end{equation}
where $\phi_{h,t}$ [MW/(m\textsuperscript{3}/s)] is the pump factor,
$\eta_{h,t} \in (0,1]$ is the stage efficiency, $\varphi_{w,s,t,b}$
[m\textsuperscript{3}/s] is the waterway flow, and
$\ell_{d,s,t,b}$ [MW] is the served electrical load.  An optional
capacity constraint bounds the flow:
$\varphi_{w,s,t,b} \leq \overline{Q}_{h,t}$.

The pump element only adds coupling rows to the LP; the waterway
retains its junction water balance contribution (Section~\ref{sec:hydro})
and the demand retains its bus power balance contribution
(Section~\ref{sec:balance}).  Reversible turbine-pump units (e.g.
pump-storage hydro) are modeled as a pair of elements: a
\texttt{Turbine} on the generation waterway and a \texttt{Pump} on the
pumping waterway.  Economic optimality (round-trip efficiency
$< 1.0$) prevents simultaneous generation and pumping without
explicit binary exclusivity constraints.

\subsection{LNG Terminal}
\label{sec:lng}

An \texttt{LngTerminal} element $\tau \in \setLNG$ models an
LNG storage tank whose inventory is coupled to gas-fired generators
through their heat rates.  The tank volume balance mirrors the battery
SoC equation (Section~\ref{sec:battery}) with delivery inflow and
boil-off:
\begin{multline}
\label{eq:lng_balance}
v_{\tau,s,t,b}^{\text{lng}}
  = v_{\tau,s,t,b-1}^{\text{lng}} \bigl(1 - \lambda_\tau^h \Delta_b\bigr)
  + \rho_\tau \Delta_b \, d_{\tau,s,t,b}^{\text{lng}} \\
  - \rho_\tau \Delta_b \, \text{vent}_{\tau,s,t,b}
  - \rho_\tau \Delta_b
      \sum_{g \in \mathcal{G}_\tau} \kappa_g^{hr} \, p_{g,s,t,b}
\end{multline}
where $v_{\tau,s,t,b}^{\text{lng}}$ [m\textsuperscript{3}] is the tank
volume, $\lambda_\tau^h = \lambda_\tau / 8760$ is the per-hour
boil-off rate, $\rho_\tau$ is a dimensionless unit-conversion factor
(default~1.0, used when delivery is in m\textsuperscript{3}/h and volume
in m\textsuperscript{3} so that $\rho_\tau \cdot \Delta_b \cdot d$ has
units m\textsuperscript{3}),
$d_{\tau,s,t,b}^{\text{lng}}$ is the delivery inflow, $\text{vent}$
is a vented-gas variable penalized at `spillway\_cost', and
$\kappa_g^{hr}$ [m\textsuperscript{3}/MWh] is the heat rate of each
linked generator $g \in \mathcal{G}_\tau$.

The last term is implemented by injecting a coefficient
$-\rho_\tau \kappa_g^{hr} \Delta_b$ onto each generator's generation
variable in the existing tank balance row.  Volume bounds and
initial/final conditions follow the generic storage template.  The
delivery schedule is fixed to a constant hourly rate
$d_{\tau,s,t,b} = D_{\tau,t} / T_t$ across each stage.

This formulation enables the optimizer to manage fuel inventory (buy
decisions and storage costs) jointly with economic dispatch, capturing
supply security constraints and seasonal LNG price dynamics.

\subsection{Unit Commitment}
\label{sec:uc}

The \texttt{Commitment} element provides thermal unit commitment with
the classical three-bin (3-bin) formulation using status~$u$,
startup~$v$, and shutdown~$w$ binary
variables~\cite{morales2013,knueven2020}.  The fundamental status
transition logic is:
\begin{align}
u_{g,s,t,p} - u_{g,s,t,p-1} - v_{g,s,t,p} + w_{g,s,t,p}
  &= 0
  \label{eq:uc_logic} \\
v_{g,s,t,p} + w_{g,s,t,p}
  &\leq 1
  \label{eq:uc_excl} \\
\underline{P}_g u_{g,s,t,p(b)}
  \leq p_{g,s,t,b}
  &\leq \overline{P}_g u_{g,s,t,p(b)}
  \label{eq:uc_bounds}
\end{align}
where $p$ indexes commitment periods (typically hourly) and $b$ indexes
dispatch blocks (sub-hourly allowed).  Binary variables live on the
commitment-period grid, enabling sub-hourly dispatch with hourly on/off
decisions.

Minimum up/down time constraints, ramp limits (while online, during
startup, and shutdown), hot/warm/cold startup cost tiers, shutdown
costs, no-load costs, and piecewise heat-rate curves are all supported.
Commitment constraints apply only to stages flagged
\texttt{chronological}~$=$ true (consecutive hourly blocks).  On
representative-period stages (duration-weighted blocks not in calendar
order), binary variables are silently relaxed to continuous.

\subsection{Simple Commitment}
\label{sec:sc}

The \texttt{SimpleCommitment} element provides a lightweight commitment
variant for generators that require minimum-output enforcement without
the full startup/shutdown tracking overhead.  It introduces a single
binary status variable per block:
\begin{equation}
u_{g,s,t,b}^{\text{sc}} \;\in\; \{0,1\}
\qquad \forall \; g \in \setUC,\; s,t,b
\label{eq:sc_var}
\end{equation}
with generation bounds:
\begin{align}
p_{g,s,t,b}
  &\leq \overline{P}_{g,t,b}\,u_{g,s,t,b}^{\text{sc}}
  \label{eq:sc_upper} \\
p_{g,s,t,b}
  &\geq \underline{P}_{g,t,b}^{d}\,u_{g,s,t,b}^{\text{sc}}
  \label{eq:sc_lower}
\end{align}
where $\underline{P}_{g,t,b}^{d}$ (\texttt{dispatch\_pmin}) is the
minimum output when dispatched, defaulting to the generator's
$\underline{P}_g$ when unset.  When the generator is offline ($u = 0$),
generation is forced to zero; when online ($u = 1$), generation must lie
in $[\underline{P}_{g}^{d}, \overline{P}_{g}]$.

Unlike the full \texttt{Commitment} element (Section~\ref{sec:uc}),
\texttt{SimpleCommitment}:
\begin{itemize}
  \item Has no startup/shutdown variables $v$, $w$ or associated costs
  \item Has no minimum up/down time constraints
  \item Has no ramp limits
  \item Has no no-load cost
  \item Has no commitment-period abstraction (binary per dispatch block)
  \item Does \emph{not} require chronological stages
\end{itemize}
This lighter formulation is sufficient for daily unit commitment in many
hydrothermal systems and for PLP-style inertia modeling where only the
online/offline status matters.  When \texttt{must\_run}~$=$ true, the
column is clamped to $u_{g,s,t,b}^{\text{sc}} = 1$ via lower bound
$\underline{u}^{\text{sc}} = 1$.  When \texttt{relax}~$=$ true or in
SDDP continuous phases, the integrality is dropped and
$u_{g,s,t,b}^{\text{sc}} \in [0,1]$.

\subsection{Emissions}
\label{sec:emission}

Emissions are modeled with three elements: a pollutant registry
(CO\textsubscript{2}, CH\textsubscript{4}, NO\textsubscript{x},
SO\textsubscript{2}, particulate matter, \dots); emission zones (a
system-wide or regional scope covering a set of pollutants with weights,
e.g.\ GWP-100 for a CO\textsubscript{2}-equivalent basket, and an
optional price and cap); and emission sources that link generators to
zones.  Sources can also be derived from fuel emission factors and heat
rates.  There is no emission variable: each source adds to its
generator's dispatch column $p_{g,s,t,b}$ the coefficient
\begin{equation}
\label{eq:emission_coef}
\alpha^{\text{em}}_{\nu,t,b} = \mu_{z,\epsilon}
  \bigl(1-\eta^{\text{cap}}_{g,\epsilon,t}\bigr)
  \bigl(f_{\nu,t} + f^{\text{up}}_{\nu,t}\bigr) \Delta_b
  \quad \text{[t/MW]},
\end{equation}
where $\mu_{z,\epsilon}$ is the zone weight, $\eta^{\text{cap}}$ the
carbon-capture fraction, and $f$, $f^{\text{up}}$ the combustion and
upstream rates [t/MWh].  A zone price $c^{\text{em}}_{z,t}$ adds
$c^{\text{em}}_{z,t}\,\alpha^{\text{em}}_{\nu,t,b}/\Delta_b$ to the
generator's cost.  A zone cap adds one row per scenario and stage,
\begin{equation}
\label{eq:emission_cap}
\sum_{b \in \setB_t} \sum_{\nu \in z}
  \alpha^{\text{em}}_{\nu,t,b}\, p_{g(\nu),s,t,b}
  - s^{\text{cap}}_{z,s,t} \leq \overline{E}_{z,t},
\end{equation}
whose dual is the endogenous price of the pollutant; the optional slack
$s^{\text{cap}}$ makes the cap soft.  For cap-and-trade, zones draw on
an allowance pool that reuses the storage balance: allocations flow in,
emissions are withdrawn, auction purchases are allowed at a price, and
unused allowances are banked across stages, so the cap becomes a
multi-stage budget.  From the bus and cap duals, GTOpt also reports
marginal emission factors per bus and hour, including the emissions
avoided or caused by storage operation.

\subsection{Multi-Carrier Chains}
\label{sec:carriers}

Hydrogen, ammonia and heat are represented by carrier-typed balance
nodes, storage elements and one-stage converters between any pair of
carriers.  A converter with efficiency $\eta_c$ withdraws its input
$x_{c,s,t,b}$ from the source node and injects $\eta_c x_{c,s,t,b}$ at
the destination node, which models electrolysers, fuel cells,
Haber--Bosch synthesis, ammonia cracking, electric boilers and
concentrated-solar power blocks.  Hydrogen, ammonia and thermal storage
reuse the storage balance of Section~\ref{sec:battery}.  Together with
the LNG terminal of Section~\ref{sec:lng}, the same LP therefore
co-optimizes the power system and its fuel and power-to-X chains.

\subsection{Synchronous Inertia Requirements}
\label{sec:inertia}

As variable renewable penetration increases, maintaining sufficient
synchronous inertia for frequency stability becomes a binding operational
constraint.  \gtopt{} models inertia requirements through two elements
mirroring the reserve-zone formulation (Section~\ref{sec:reserve}).

An \texttt{InertiaZone} $z \in \setIZ$ specifies a minimum inertia
requirement $\overline{R}_{z,t,b}^{H}$ [MWs]:
\begin{equation}
\label{eq:inertia_req}
\sum_{p \in \mathcal{P}_z} \Phi_{p,t} \, r_{p,s,t,b}^{H}
  \geq \overline{R}_{z,t,b}^{H}
\end{equation}
An \texttt{InertiaProvision} $p \in \setIP$ links a generator to
one or more zones.  The provision variable
$r_{p,s,t,b}^{H} \in [0, \overline{R}_{p,t,b}^{H}]$ [MW] is bounded
by a downward coupling row:
\begin{equation}
\label{eq:inertia_coupling}
p_{g,s,t,b} \geq r_{p,s,t,b}^{H}
\end{equation}
The conversion factor from MW dispatch to MWs of inertia is:
\begin{equation}
\Phi_{p,t} = \frac{H_p \cdot S_p}{\underline{P}_g}
\end{equation}
where $H_p$ [s] is the machine inertia constant and $S_p$ [MVA] is the
rated apparent power.  A user-specified \texttt{provision\_factor} may
override this formula.  Shortage is penalized at the zone's
\texttt{cost} [\$/MWs], enabling the optimizer to trade off committed
generation against inertia value.

\subsection{Water Rights and Irrigation Agreements}
\label{sec:rights}

Hydro-dominated power systems often operate under legally-binding
irrigation agreements that partition turbine discharge into distinct
rights categories---irrigation, electrical, mixed---with
volume-dependent allocations that vary with reservoir level.  In Chile,
the Laja Agreement (1958, updated 2017) and the Maule Convention (1947)
govern two of the largest hydro cascades in the national interconnected
system, imposing seasonal flow obligations and multi-annual volume
accounting constraints on the optimizer.

Traditional tools such as PLP encode each agreement as a monolithic
Fortran module, tightly coupling the rights logic to the physical
topology.  \gtopt{} instead introduces two generic, data-driven
entities---\emph{FlowRight} and \emph{VolumeRight}---together with
the flexible \emph{UserConstraint} mechanism, forming a
\emph{rights-domain} layer parallel to the physical hydro topology.
This separation allows arbitrary irrigation agreements to be specified
declaratively in JSON and parameterized AMPL-inspired scripts without
modifying the solver code.

\subsubsection{Flow Right}

A flow right $i \in \setFR$ represents a flow-based water right
[m$^3$/s] that creates one LP column per block.  Two operating modes
are supported:
\begin{itemize}
  \item \emph{Fixed mode}: the extraction flow is fixed at a scheduled
    discharge value $\bar{\xi}_{i,s,t,b}$ (lower bound equals upper
    bound).
  \item \emph{Variable mode}: when a maximum flow $\overline{F}_i > 0$
    is specified, the optimizer decides the flow
    $\xi_{i,s,t,b} \in [0, \overline{F}_i]$.
\end{itemize}

When $\texttt{consumptive} = \texttt{true}$, the flow right subtracts
from the physical junction balance, coupling the rights domain to the
physical hydro topology.  An optional deficit variable with penalty
cost~$c_i^f$ ensures feasibility when the right cannot be fully served.

A stage-average hourly flow variable $\bar{\xi}^h_i$ is optionally
created to aggregate block-level flows into the duration-weighted
average required by the volume accumulation equations:
\begin{equation}
\label{eq:fr_avg}
\bar{\xi}^h_{i,s,t} = \sum_{b \in \setB_t}
  \frac{\Delta_b}{T_t} \, \xi_{i,s,t,b}
\end{equation}

\subsubsection{Volume Right}

A volume right $o \in \setVR$ is modeled as a virtual reservoir
(following the same \texttt{StorageLP} pattern as batteries) that
tracks accumulated volume [hm$^3$] across stages.  The volume balance
is:
\begin{equation}
\label{eq:vr_balance}
\nu_{o,s,t} = \nu_{o,s,t-1}
  - \alpha_o \sum_{b \in \setB_t} \Delta_b \, \xi^\text{in}_{o,s,t,b}
\end{equation}
where $\nu_{o,s,t}$ is the remaining volume of right~$o$ at the end of
stage~$t$, $\xi^\text{in}_{o,s,t,b}$ is the input extraction rate
[m$^3$/s], and $\alpha_o = 0.0036$~[hm$^3$/(m$^3$/s$\cdot$h)] is the
flow-to-volume conversion factor.

Volume rights support inter-stage state-variable coupling for SDDP,
optional demand satisfaction constraints
($\alpha_o \sum_b \Delta_b \xi^\text{in}_b + q^\text{fail} \geq D_o$),
and an automatic seasonal reset mechanism: when a
\texttt{reset\_month} is configured, the accumulated volume is reset to
the initial value at the start of each hydrological year.

\subsubsection{User Constraints and Parameters}
\label{sec:user_constraints}

The partition logic that links flow rights together---and the
proportional allocation rules of the Maule Convention---is expressed
through \emph{UserConstraints}: user-defined linear constraints written
in an AMPL-inspired pseudo-script.  Each constraint is a named linear
(in)equality over LP variables of any system element (generators,
demands, flow rights, volume rights, etc.) and can optionally restrict
its domain to specific scenarios, stages, or blocks.

For the Laja agreement, the total turbine flow (``gasto turbinado'')
must equal the sum of all partitioned extraction categories---irrigation
rights (``derechos de riego''), electrical rights (``derechos
el\'ectricos''), mixed rights (``derechos mixtos''), and anticipated
discharge (``gasto anticipado''):
\begin{equation}
\label{eq:uc_partition}
\xi^\text{gen} = \xi^\text{irr} + \xi^\text{elec}
  + \xi^\text{mix} + \xi^\text{antic}
\end{equation}
In the \texttt{.pampl} template this constraint reads:
{\footnotesize
\begin{verbatim}
constraint laja_particion_derechos
  "Flow partition":
 flow_right('laja_q_turbinado').flow =
   flow_right('laja_der_riego').flow
 + flow_right('laja_der_electrico').flow
 + flow_right('laja_der_mixto').flow
 + flow_right('laja_gasto_anticipado').flow;
\end{verbatim}
}

Entity names follow the PLP Spanish naming conventions (e.g.,
\texttt{laja\_der\_riego} for PLP's \texttt{IQDR}, ``Derecho de
Riego'') to maintain traceability with the reference Fortran
implementation.  This single UserConstraint replaces the hardcoded
balance node that traditional tools implement as a dedicated data
structure.

When a FlowRight has a positive \texttt{fail\_cost} and a nonzero
\texttt{discharge}, a deficit coupling constraint is automatically
generated:
\begin{equation}
\label{eq:deficit_coupling}
\xi_{i,b} + \delta_{i,b} \geq d_{i,b}
\end{equation}
where $\delta_{i,b}$ is the deficit variable penalized in the
objective.  This ensures the optimizer only uses the deficit
variable when the flow is genuinely constrained by junction
balance or dynamic bounds.

For the Maule Convention, the ordinary reserve zone enforces a
proportional split between electric ($p_e$) and irrigation ($p_r$)
extraction from the same shared flow:
\begin{align}
\label{eq:uc_maule_pct}
\xi^\text{elec}_\text{ord} &\leq
  p_e \bigl(\xi^\text{elec}_\text{ord} + \xi^\text{irr}_\text{ord}\bigr) \\
\xi^\text{irr}_\text{ord}  &\leq
  p_r \bigl(\xi^\text{elec}_\text{ord} + \xi^\text{irr}_\text{ord}\bigr)
\end{align}
where $p_e + p_r = 1$.  In \gtopt{}, these constraints reference the
Maule FlowRights using Spanish naming
(\texttt{maule\_gasto\_ordinario\_elec} for PLP's \texttt{IQMOE},
\texttt{maule\_gasto\_ordinario\_riego} for \texttt{IQMOR}),
replacing the hardcoded Fortran proportional-split logic in PLP's
\texttt{genpdmaulen.f}.

Named scalar or monthly-indexed \emph{UserParams} ($p_e$, $p_r$, cost
arrays, seasonal demand curves, etc.) are defined in the same script
file and resolved automatically from the stage calendar month, enabling
seasonally-varying constraints without per-stage hardcoding.
Each agreement is maintained in a separate \texttt{.pampl} file
(e.g., \texttt{laja\_agreement.pampl}, \texttt{maule\_agreement.pampl})
and loaded via the \texttt{user\_constraint\_files} array field in the
system JSON, keeping each agreement self-contained.
The \texttt{plp2gtopt} converter generates these templates automatically
from the parsed Fortran agreement parameters.

\subsubsection{Volume-Dependent Dynamic Bounds}

A key feature of irrigation agreements is that the available rights
depend on the current reservoir volume.  For example, the Laja agreement
defines four volume zones that progressively allocate rights as the
reservoir fills:
\begin{equation}
\label{eq:bound_rule}
\overline{F}_i(V) = \min\!\Bigl(
  c_k + m_k (V - V_k),\; \overline{F}_i^\text{cap}
\Bigr)
\end{equation}
where $(V_k, m_k, c_k)$ are breakpoints of a piecewise-linear function
and $V$ is the current volume of the source reservoir.  At each SDDP
iteration, the bound is re-evaluated from the forward-pass reservoir
volume and applied via \texttt{set\_col\_upp} on the corresponding flow
or extraction column.

Table~\ref{tab:laja_zones} illustrates the Laja volume zone
allocation.

\begin{table}[t]
\centering
\caption{Laja Agreement volume zone allocation factors.}
\label{tab:laja_zones}
\small
\begin{tabular}{@{}lrccc@{}}
\toprule
\textbf{Zone} & \textbf{Width} & \textbf{Irrig.}
  & \textbf{Elec.} & \textbf{Mixed} \\
  & [hm$^3$] & factor & factor & factor \\
\midrule
Base       & ---     & 570    & 0      & 30   \\
Zone~1     & 1\,200  & 0.00   & 0.05   & 1.00 \\
Zone~2     & 170     & 0.40   & 0.05   & 0.00 \\
Zone~3     & 530     & 0.40   & 0.40   & 0.00 \\
Zone~4     & 3\,682  & 0.25   & 0.65   & 0.00 \\
\bottomrule
\end{tabular}
\\[4pt]
{\footnotesize Base values in m$^3$/s; zone factors in
m$^3$/s per hm$^3$.  Caps: 5\,000 / 1\,200 / 30~m$^3$/s.}
\end{table}

\subsubsection{Design Advantages}

The rights-domain abstraction provides three advantages over the
monolithic approach used in traditional tools: (i)~new agreements are
specified declaratively in JSON and AMPL-inspired scripts without
modifying the solver code; (ii)~the same two generic entities
(FlowRight and VolumeRight) together with UserConstraints express both
the Laja and Maule agreements despite their very different structures
(4 volume zones vs.\ 3 reservoir portions, multi-annual vs.\ seasonal
reset, exact partition vs.\ proportional split); (iii)~volume-dependent
dynamic bounds integrate naturally with the SDDP forward-pass update
mechanism, reusing the same piecewise-linear evaluation used for
turbine efficiency.  The \texttt{plp2gtopt} converter automatically
transforms PLP's Fortran agreement parameters into the equivalent
\gtopt{} JSON configuration and \texttt{.pampl} template with verified
numerical equivalence.

\subsection{Constraint Summary}

Table~\ref{tab:constraints} lists all constraint families in the
formulation.

\begin{table}[t]
\centering
\caption{Complete constraint summary.}
\label{tab:constraints}
\small
\begin{tabular}{@{}clll@{}}
\toprule
\textbf{\#} & \textbf{Constraint} & \textbf{Eq.} & $\boldsymbol{\forall}$ \\
\midrule
C1  & Bus power balance             & \eqref{eq:balance}          & $n,s,t,b$ \\
C2  & Demand balance                & \eqref{eq:demand_balance}   & $d,s,t,b$ \\
C3  & Kirchhoff voltage law         & \eqref{eq:kirchhoff}        & $\ell,s,t,b$ \\
C4  & Generator output bounds       & \eqref{eq:pgen}             & $g,s,t,b$ \\
C5  & Capacity expansion            & \eqref{eq:capacity}         & $g,t$ \\
C6  & Generator profile             & \eqref{eq:profile}          & $g,s,t,b$ \\
C7  & Demand min.\ energy           & \eqref{eq:min_energy}       & $d,s,t$ \\
C8  & Line capacity                 & ---                         & $\ell,s,t,b$ \\
C9  & Battery SoC                   & \eqref{eq:soc_balance}      & $e,s,t,b$ \\
C10 & Converter coupling            & \eqref{eq:conv_discharge}   & $v,s,t,b$ \\
C11 & Reserve requirements          & \eqref{eq:reserve_up}       & $z,s,t,b$ \\
C12 & Reserve--generator coupling   & ---                         & $p,s,t,b$ \\
C13 & Hydro water balance           & \eqref{eq:hydro_balance}    & $j,s,t,b$ \\
C14 & Turbine conversion            & \eqref{eq:turbine}          & $h,s,t,b$ \\
C15 & Transformer / PST             & \eqref{eq:transformer}      & $\ell,s,t,b$ \\
C16 & Flow right averaging          & \eqref{eq:fr_avg}           & $i,s,t$ \\
C17 & Volume right balance          & \eqref{eq:vr_balance}       & $o,s,t$ \\
C18 & User-defined constraint       & \eqref{eq:uc_partition}     & user-defined \\
C19 & Pump power--flow coupling     & \eqref{eq:pump}             & $h,s,t,b$ \\
C20 & LNG tank volume balance       & \eqref{eq:lng_balance}      & $\tau,s,t,b$ \\
C21 & Unit commitment logic (3-bin) & \eqref{eq:uc_logic}         & $g,s,t,p$ \\
C22 & UC generation bounds          & \eqref{eq:uc_bounds}        & $g,s,t,b$ \\
C23 & Simple commitment variable    & \eqref{eq:sc_var}           & $g,s,t,b$ \\
C24 & SC generation bounds (upper)  & \eqref{eq:sc_upper}         & $g,s,t,b$ \\
C25 & SC generation bounds (lower)  & \eqref{eq:sc_lower}         & $g,s,t,b$ \\
C26 & Inertia zone requirement      & \eqref{eq:inertia_req}      & $z,s,t,b$ \\
C27 & Inertia provision coupling    & \eqref{eq:inertia_coupling} & $p,s,t,b$ \\
C28 & Loss tangents (signed flow)   & \eqref{eq:loss_tangent}     & $\ell,s,t,b,k$ \\
C29 & Emission cap                  & \eqref{eq:emission_cap}     & $z,s,t$ \\
C30 & Gas node balance              & \eqref{eq:gas_balance}      & $i,s,t,b$ \\
C31 & Gas supply stage volume       & \eqref{eq:gas_volume}       & $k,s,t$ \\
C32 & Weymouth tangents             & \eqref{eq:gas_weymouth}     & $a,s,t,b,m$ \\
C33 & Line pack                     & \eqref{eq:gas_linepack_bal} & $a,s,t,b$ \\
C34 & Compressor ratio              & \eqref{eq:gas_ratio}        & $c,s,t,b$ \\
C35 & Compressor fuel plane         & \eqref{eq:gas_fuel_plane}   & $c,s,t,b$ \\
C36 & Plant pressure gate / derate  & \eqref{eq:gas_gate}         & $g,s,t,b$ \\
\bottomrule
\end{tabular}
\end{table}
